\subsection*{Calculations with centered prime sums}

\teamanchor{proof.centered-ramp-count}%
Let \(\Lambda\) be the von Mangoldt function and \(v_+=\max(v,0)\).
Take complex coefficients \(b(n)\) with \(|b(n)|\ll\Lambda(n)\).
For \(y\ge1\) and constants \(a,c\), put
\[
 A_b(y)=\sum_{n\le y}b(n),\qquad H_b(y)=A_b(y)-ay,\qquad
 \mathcal R_b(y)=\sum_{n\ge1}b(n)(1/n-1/y)_+-a\log y-c .
\]
Write \(D_b(s)=\sum_{n\ge1}b(n)n^{-s}\) for \(\Re s>1\).
Integrating the step function \(A_b\) gives
\begin{equation}\label{eq:common-ramp-count}
 \mathcal R_b(y)+c=\int_1^y H_b(t)t^{-2}\,dt,\qquad
 \int_1^\infty H_b(t)t^{-s-1}\,dt
   =\frac{D_b(s)}s-\frac{a}{s-1}\quad(\Re s>1).
\end{equation}
Indeed, a term at \(n\) contributes the integral over \(t\ge n\).
For the infinite integral, the absolute sum of those integrals at
\(\sigma=\Re s>1\) is \(\sigma^{-1}\sum_n |b(n)|n^{-\sigma}<\infty\).

\teamanchor{proof.mellin-ramp-kernel}%
Suppose \(D_b\) continues meromorphically to \(\Re s>0\) and
\(C_b(s)=D_b(s)-as/(s-1)\) is holomorphic at \(1\).
Choose \(c=C_b(1)\). For \(\Re s>1\), the Mellin integral of
\(y\mathcal R_b(y)\) converges absolutely and satisfies
\begin{equation}\label{eq:common-ramp-mellin}
 \Phi_b(s)=\int_1^\infty y\mathcal R_b(y)y^{-s-1}\,dy
 =\frac{D_b(s)}{s(s-1)}-\frac{a}{(s-1)^2}-\frac{c}{s-1}
 =\frac{C_b(s)/s-c}{s-1}.
\end{equation}
The integral of one ramp is
\(\int_n^\infty (y/n-1)y^{-s-1}\,dy=n^{-s}/[s(s-1)]\).
At \(\sigma=\Re s>1\) the sum of the absolute integrals is
\(\sum_n |b(n)|n^{-\sigma}/[\sigma(\sigma-1)]<\infty\).
Integration of \(y\log y\) and \(y\) gives the remaining two terms.
The last expression makes \(1\) removable. At a simple pole
\(\rho\ne1\) of \(D_b\) in \(\Re s>0\), the residue of \(\Phi_b\) is
\(\operatorname{Res}_{s=\rho}D_b(s)/[\rho(\rho-1)]\).

If \(\int_1^\infty |H_b(t)|t^{-2}\,dt<\infty\), take the limit
along real \(s\downarrow1\) in \eqref{eq:common-ramp-count}.
By dominated convergence its value is \(C_b(1)=c\). Consequently
\begin{equation}\label{eq:common-ramp-tail}
 \mathcal R_b(y)=-\int_y^\infty H_b(t)t^{-2}\,dt .
\end{equation}

\teamanchor{proof.centered-ramp-maximum}%
For real \(b(n),a,c\), fix \(0<\alpha<1\) and \(\kappa>0\).
The continuous function \(D(y)=\mathcal R_b(y)-\kappa y^{\alpha-1}\)
attains a maximum on every nonempty compact set of scales in \([1,\infty)\).
The maximizers form a compact set, so there is a least one \(y_*\).
If \(\mathcal R_b(y_*)\le \kappa y_*^{\alpha-1}\), then
\(\mathcal R_b(y)\le \kappa y^{\alpha-1}\) at every scale in the set.
The constant \(c\) does not affect the choice.

On a real interval \([Y,Z]\), \eqref{eq:common-ramp-count} gives
\begin{equation}\label{eq:common-ramp-maximum}
 D(y)-D(Y)=\int_Y^y
      \frac{A_b(t)-at+(1-\alpha)\kappa t^\alpha}{t^2}\,dt
       \qquad (Y\le y\le Z).
\end{equation}
If \(a>\alpha(1-\alpha)\kappa Y^{\alpha-1}\), the function
\(at-(1-\alpha)\kappa t^\alpha\) is strictly increasing on this interval.
Between consecutive changes of \(A_b\) it follows from the integrand
that \(D\) has at most one internal critical point.
The first piece uses \(A_b(Y)\), and each interior support point
adds its coefficient to the count.
It suffices to compare the endpoints, the support points of \(b\)
in \((Y,Z)\), and the interior point where
\(at-(1-\alpha)\kappa t^\alpha=A_b(t)\) in each open piece if it exists.
For a finite set of scales one can compare its values directly.

\subsection*{The signed remainder as a weighted prime sum}

\teamanchor{proof.weighted-prime-series}%
Keep \(M=18480\) and a fixed prime \(d\equiv9239\pmod M\). Write
\(\chi_d(n)=(n/d)\) for the Legendre symbol. For \(Q\in\{M,1\}\) and
real \(y\ge1\), define
\begin{align*}
 \lambda_{d,Q}(n)&=\mathbf{1}_{n\equiv1\bmod Q}\Lambda(n)(1-\chi_d(n)),\\
 S_{d,Q}(y)&=\sum_{n\le y}\lambda_{d,Q}(n)\left(\frac1n-\frac1y\right),\\
 R_{d,Q}(y)&=S_{d,Q}(y)-\frac{\log y}{\varphi(Q)}-c_{d,Q}.
\end{align*}
At \(Q=1\) the progression condition is automatic, with \(\varphi(1)=1\),
and the weight is still \(\Lambda(n)(1-\chi_d(n))\).

For any Dirichlet character \(\xi\), let \(\widetilde L_\xi(s)=L(s,\xi)\) for
nonprincipal \(\xi\) and let \(\widetilde L_\xi(s)=(s-1)L(s,\xi)\) for
principal \(\xi\), with the value at \(1\) supplied by continuation. The
twist \(\xi\chi_d\) is a character modulo \(Qd\). Euler products are taken
at the declared moduli. The regularized functions are holomorphic on
\(\Re s>0\) and nonzero on \(\Re s\ge1\), as shown in the
\hyperlink{proof.dirichlet-zero-free-line}{boundary calculation}.
Put
\begin{equation}\label{eq:weighted-centered-log-derivative}
 G_{d,Q}(s)=\frac{1}{\varphi(Q)}
 \left(\sum_{\xi\bmod Q}\left(
 \frac{\widetilde L'_{\xi\chi_d}(s)}{\widetilde L_{\xi\chi_d}(s)}
 -\frac{\widetilde L'_\xi(s)}{\widetilde L_\xi(s)}\right)-1\right),
 \qquad c_{d,Q}=\Re G_{d,Q}(1).
\end{equation}
Since \((d,Q)=1\), the Chinese remainder theorem gives a unit modulo
\(Qd\) that is \(1\) modulo \(Q\) and on which \(\chi_d\) is \(-1\).
Hence every twist is nonprincipal. Let
\(F_{d,Q}(s)=\sum_{n\ge1}\lambda_{d,Q}(n)n^{-s}\) for \(\Re s>1\).
For every \(n\), character orthogonality gives
\(\varphi(Q)\mathbf 1_{n\equiv1\bmod Q}(1-\chi_d(n))
=\sum_{\xi\bmod Q}(\xi(n)-(\xi\chi_d)(n))\).
This includes multiples of \(d\), at which the twisted term vanishes.
The differentiated Euler products \eqref{eq:common-character-series}
therefore give
\begin{equation}\label{eq:weighted-series}
 F_{d,Q}(s)=\frac{s}{\varphi(Q)(s-1)}+G_{d,Q}(s).
\end{equation}
Only the untwisted principal character contributes the pole term.
The coefficients of \(F_{d,Q}\) are real, so \(G_{d,Q}\) is real for
real \(s>1\) and, by continuation, at \(1\). Thus \(c_{d,Q}=G_{d,Q}(1)\).
At \(Q=1\) there is just the untwisted function \((s-1)\zeta(s)\) and
its twist \(L(s,\chi_d)\). The Laurent expansion of \(\zeta\) at \(1\) gives
\(\widetilde L'_1(1)/\widetilde L_1(1)=\gamma\), so
\(c_{d,1}=-1-\ell_d\) for
\(\ell_d=\gamma-\Re(L'(1,\chi_d)/L(1,\chi_d))\).\footnote{\href{https://github.com/lsernau/grh-lean/blob/72fead102f663d5c7a407964a13dc0cf9d977207/grh-compact/src/New/Endgame/PrimeError.lean\#L23-L65}{Compact rewrite, definition of the error and evaluation at modulus one}.}

\teamanchor{proof.weighted-prime-identity}%
Take the \hyperlink{proof.weighted-data}{same datum} used for the row
comparison, with set \(\mathcal T\) of top places, set \(\mathcal J\) of
thick places, \(T=\#\mathcal T\), \(J=\#\mathcal J\) and
\(f=\varphi(M)T\). Recall that top places are the class \(1\) in the
reduction of the place group from conductor \(N_1\) to \(M\), and thick
places belong to the class \(-1\). If \(n\) is a unit modulo \(M\), it is a
unit modulo \(N_1\) because \(N_1/M\) has only the prime factors \(2,3\),
which divide \(M\). Then \(\tau\mapsto[-n]\tau\) sends the top fiber
bijectively onto the class \(-n\). For the other \(n\) the place \([-n]\)
is the identity. In either case the number of top places for which
\([-n]\tau\) is thick is \(J\mathbf{1}_{n\equiv1\bmod M}\).

In \eqref{eq:weighted-remainder}, the nominal weight \(w\) equals \(a/x\)
at a thick place and \(a/x_1\) elsewhere. Both scales are at most \(a\),
since \(a\ge xb\)
and \(\log x_1\le\log b\) for these data. In the sum
\(\operatorname{Ask}(a,\tau)\), the positive part of
\(a/n-w([-n]\tau)\), divided by \(a\), is therefore the kernel of
\(S_{d,1}(x)\) or of \(S_{d,1}(x_1)\) according to the translated
place. Counting the two kernels over the top places gives
\[
 \frac1a\sum_{\tau\in\mathcal T}\operatorname{Ask}(a,\tau)
 =J\bigl(S_{d,M}(x)-S_{d,M}(x_1)\bigr)+T S_{d,1}(x_1).
\]
The top sum of the main ask is evaluated using the same weights.
Specifically \(\bar w=f^{-1}\sum_\sigma\log w(\sigma)\) is
\(\log a-\log x_1-(J/f)(\log x-\log x_1)\), so
\[
 \frac{T}{a}\operatorname{MainAsk}
 =\frac{J}{\varphi(M)}(\log x-\log x_1)+T(\log x_1+c_{d,1}).
\]
Here the modulus-one evaluation \(c_{d,1}=-1-\ell_d\) is the one from
\hyperlink{proof.weighted-prime-series}{the series calculation}.
The two occurrences of \(c_{d,M}\) in the progression difference cancel.
Subtracting the main ask from each term of the top sum gives the identity
for the remainder in the \hyperlink{proof.weighted-height}{height inequality},
\begin{equation}\label{eq:weighted-prime-identity}
 \operatorname{rem}=a\{J(R_{d,M}(x)-R_{d,M}(x_1))+T R_{d,1}(x_1)\}.
\end{equation}
This is the counting argument in the compact
\href{https://github.com/lsernau/grh-lean/blob/72fead102f663d5c7a407964a13dc0cf9d977207/grh-compact/src/New/Endgame/ShortAndEnough.lean\#L84-L112}{remainder identity}.

\subsection*{The error at the supplied second scale}

\teamanchor{proof.weighted-second-scale}%
\teamanchor{proof.fixed-modulus-pnt}%
For each \(A>0\), Koukoulopoulos's progression theorem gives
uniformly for \(t\ge3\), \(u\le(\log t)^A\) and \((r,u)=1\)
\[
 \sum_{n\le t,\ n\equiv r\bmod u}\Lambda(n)
 =\frac{t}{\varphi(u)}
  +O_A\!\left(t\exp\!\left(-c_A(\log t)^{3/5}
                             (\log\log t)^{-1/5}\right)\right).
\]
The positive constant \(c_A\) is ineffective.
\cite[Theorem~1.1, arXiv v3, p.~2]{f1koukoulopoulos}
Take \(A=1\). Each fixed \(u\ge2\) satisfies the modulus bound
eventually, and at this fixed modulus the error is
\(O_{u,B}(t/(\log t)^{B+1})\) for every \(B>0\).
For \(u=1\) take the odd class modulo \(2\) and restore the powers of
\(2\), which have total weight \(O(\log t)\). Since
\(\varphi(1)=\varphi(2)=1\), for every fixed \(u\ge1\), reduced class
\(r\bmod u\) and \(B>0\), as real \(t\to\infty\),
\begin{equation}\label{eq:weighted-fixed-pnt}
 \sum_{n\le t,\ n\equiv r\bmod u}\Lambda(n)
 =\frac{t}{\varphi(u)}+O_{u,B}\!\left(\frac{t}{(\log t)^{B+1}}\right).
\end{equation}
For the weighted estimate the moduli \(1,d,M,Md\) are held fixed as
\(t\) varies. The selection of the twist prime below uses only
\(u=M\). The applications need only eventual estimates at those moduli.

\teamanchor{proof.dirichlet-zero-free-line}%
For any character \(\xi\bmod u\), let \(\delta_\xi\) be \(1\) if \(\xi\)
is principal and \(0\) otherwise. The Euler product, differentiated
in \(\Re s>1\), gives
\begin{equation}\label{eq:common-character-series}
 \sum_{n\ge1}\frac{\Lambda(n)\xi(n)}{n^s}
  =-\frac{L'(s,\xi)}{L(s,\xi)}
  =\delta_\xi\frac{s}{s-1}
       -\frac{\widetilde L'_\xi(s)}{\widetilde L_\xi(s)}-\delta_\xi .
\end{equation}
\cite[(25.2.11), (25.15.1)--(25.15.4)]{dlmf}
The product gives nonvanishing for \(\Re s>1\). The nonprincipal
\(L\)-functions are entire, and the factor \(s-1\) removes the
principal pole. Summing \eqref{eq:weighted-fixed-pnt} against \(\xi\)
gives \(H_b(t)=O_{u,B}(t/(\log t)^{B+1})\) for
\(b(n)=\Lambda(n)\xi(n)\), \(a=\delta_\xi\) in
\eqref{eq:common-ramp-count}. Therefore
\(-\widetilde L'_\xi(s)/\widetilde L_\xi(s)-\delta_\xi
=s\int_1^\infty H_b(t)t^{-s-1}\,dt\) there. The integral has a finite
limit along each horizontal approach from the right to \(\Re s=1\).
A zero of \(\widetilde L_\xi\) at any such point would give a pole
of its logarithmic derivative. This proves nonvanishing of
\(\widetilde L_\xi\) on \(\Re s\ge1\) for every fixed modulus,
including \(u=1\).

For \(Q\in\{M,1\}\), write \(A_{d,Q}(t)=\sum_{n\le t}\lambda_{d,Q}(n)\)
and \(H_{d,Q}(t)=A_{d,Q}(t)-t/\varphi(Q)\).
Apply \eqref{eq:weighted-fixed-pnt} to the reduced classes modulo \(Qd\)
that are \(1\) modulo \(Q\). By the Chinese remainder theorem the weights
\(1-\chi_d\) on these classes sum to \(d-1\), and
\(\varphi(Qd)=(d-1)\varphi(Q)\). Among the remaining terms, only powers
of \(d\) can have \(\Lambda(n)\ne0\) and \(n\equiv1\bmod Q\), and their
total weight is \(O(\log t)\). Thus for every \(B>0\),
\[
 H_{d,Q}(t)=O_{Q,d,B}\!\left(t(\log t)^{-B-1}\right).
\]
With \(b=\lambda_{d,Q}\), \(a=1/\varphi(Q)\), the series
\eqref{eq:weighted-series} gives \(C_b=G_{d,Q}\). Then
\eqref{eq:common-ramp-count} reads
\begin{equation}\label{eq:weighted-centered-count}
 \int_1^\infty H_{d,Q}(t)t^{-s-1}\,dt
 =\frac{F_{d,Q}(s)}{s}-\frac{1}{\varphi(Q)(s-1)}
 =\frac{G_{d,Q}(s)}{s} \qquad (\Re s>1).
\end{equation}
Since \(c_{d,Q}=G_{d,Q}(1)\) by the series calculation, this is the
centering in \eqref{eq:common-ramp-tail}. For all sufficiently large real \(y\),
\begin{equation}\label{eq:weighted-fixed-remainder}
 R_{d,Q}(y)=-\int_y^\infty H_{d,Q}(t)t^{-2}\,dt
 =O_{Q,d,B}\big((\log y)^{-B}\big).
\end{equation}
This verifies the centering in the compact identity at both moduli.\footnote{The
\href{https://github.com/lsernau/grh-lean/blob/72fead102f663d5c7a407964a13dc0cf9d977207/grh/src/New/Reduction/Zeros/PrimeNumberTheorem.lean\#L211-L237}{readable rewrite's fixed-modulus estimate}
uses the same centering.}

Fix \(0<\varepsilon<1/100\) as well as \(d\). Let \(x\) be a sufficiently
large first scale carrying the matching datum and row comparison.
The selected argument uses the \hyperlink{proof.weighted-square-selection}{test squares}
specified below. The original auxiliary scale satisfies
\(\log x_1=x^{\varepsilon/6}\), and \(\log b=x^\varepsilon\).
Choose \(B>12/\varepsilon\) in \eqref{eq:weighted-fixed-remainder}.
The implied constants at the two fixed moduli are then absorbed by the
exponent margin, so \(\lvert R_{d,M}(x_1)\rvert\) and
\(\lvert R_{d,1}(x_1)\rvert\) are at most \(x^{-2}\) for all sufficiently
large \(x\). Divide \(m\operatorname{rem}\le fa(\log b)^5\) by \(maJ>0\).
Using \eqref{eq:weighted-prime-identity}, \(f=\varphi(M)T\),
\(T/J\le\log b\), and \(m\ge\sqrt x\) gives
\begin{align*}
 R_{d,M}(x)
 &\le \frac{f(\log b)^5}{mJ}
       +R_{d,M}(x_1)-\frac{T}{J}R_{d,1}(x_1)\\
 &\le \varphi(M)x^{-1/2+6\varepsilon}+(1+x^\varepsilon)x^{-2}.
\end{align*}
Since \(x^{2\varepsilon}\) eventually exceeds \(\varphi(M)+2\), this is
\begin{equation}\label{eq:weighted-square-bound}
 R_{d,M}(x)\le x^{-1/2+8\varepsilon}
 \qquad\text{at every such sufficiently large }x.
\end{equation}
The fixed-modulus theorem controls the errors at the supplied second scale.
The saving at the first scale comes from the comparison on the common rows.

\subsection*{One twist prime for the entire character family}

\teamanchor{proof.weighted-twist-density}%
For an integer \(k\ge13\), put \(\beta_k=1/2+1/k\) and
\(\varepsilon_k=1/(16k)\). We need a single prime \(d_k\) in the class
\(a_0=M/2-1=9239\pmod M\) with
\begin{equation}\label{eq:weighted-twist-free}
 \widetilde L_{\xi\chi_{d_k}}(s)\ne0
 \qquad\bigl(\xi\bmod M,\ \Re s\ge\beta_k,\ |\Im s|\le k\bigr).
\end{equation}
The class is a unit because \(a_0^2\equiv1\pmod M\). Its primes satisfy
\(d_k\equiv7\pmod8\), \(M\mid2(d_k+1)\) and \(M\nmid d_k+1\), as
required by the arithmetic datum.

\paragraph{Conductors and Euler factors.}
\teamanchor{proof.coprime-character-product}%
The product of primitive characters of pairwise coprime conductors
is primitive of conductor equal to their product. In fact, the
Chinese remainder theorem lets us restrict the character to the
units of one factor with all other coordinates equal to \(1\).
If the product were induced from a proper divisor of the product
modulus, this restriction would factor through a proper divisor
of some factor conductor, contrary to primitivity.

\teamanchor{proof.dirichlet-change-of-conductor}%
If \(\xi\) modulo \(u\) is induced by \(\xi^*\) of conductor \(f\),
their Euler products give the meromorphic identity
\begin{equation}\label{eq:common-character-conductor}
 L(s,\xi)=L(s,\xi^*)\prod_{\substack{p\mid u\\p\nmid f}}(1-\xi^*(p)p^{-s}).
\end{equation}
\cite[(25.15.4)]{dlmf}
These finite factors are nonzero on \(\Re s>0\), since
\(|\xi^*(p)p^{-s}|<1\) there.

\paragraph{The candidate twists.}
\teamanchor{proof.weighted-twist-conductor}%
For a fixed character \(\xi\bmod M\), let \(\xi^*\) be its primitive
associate of conductor \(f_\xi\mid M\). At a candidate prime \(d\) in the
class \(a_0\), the character \(\psi_d=\xi^*\chi_d\) is primitive of
conductor \(f_\xi d\) by the
\hyperlink{proof.coprime-character-product}{coprime conductor rule},
and is nonprincipal because \(d>1\). In particular
\(\widetilde L_{\xi\chi_d}(s)=L(s,\xi\chi_d)\) at its declared modulus
\(Md\), so \eqref{eq:common-character-conductor} gives
\[
 \widetilde L_{\xi\chi_d}(s)=L(s,\psi_d)
      \prod_{\substack{p\mid M\\p\nmid f_\xi}}(1-\psi_d(p)p^{-s}).
\]
A failure of
\eqref{eq:weighted-twist-free} for this candidate therefore supplies a zero
of \(L(s,\psi_d)\) in \(\beta_k\le\Re s\le1\), \(|\Im s|\le k\), by
nonvanishing for \(\Re s>1\).

\teamanchor{proof.weighted-twist-exceptions}%
\teamanchor{proof.weighted-nonreal-twist-density}%
For each fixed \(\xi\bmod M\), write \(\Delta=\beta_k-1/2=1/k\).
Let \(\mathcal P(X)\) be the candidate primes \(d\le X\) in the class \(a_0\),
and let \(\mathcal B_\xi(X)\) consist of those for which
\(L(s,\psi_d)\) has a zero in the box
\[
 \beta_k\le\Re s\le1,\qquad |\Im s|\le k .
\]
We will bound \(\#\mathcal B_\xi(X)\) by \(O_{M,k}(X^{1-\Delta^2/16})\).
The compact and readable rewrites obtain a density bound by detecting
zeros with Gaussian theta integrals. The argument below uses two finite
Dirichlet polynomials and the quadratic large sieve.
\footnote{See the \href{https://github.com/lsernau/grh-lean/blob/72fead102f663d5c7a407964a13dc0cf9d977207/grh-compact/src/Lib/Analytic/QuadraticTwist/ZeroFreeTwist.lean\#L144-L186}{compact density bound} and the
\href{https://github.com/lsernau/grh-lean/blob/72fead102f663d5c7a407964a13dc0cf9d977207/grh/src/Lib/Analytic/QuadraticTwist.lean\#L679-L722}{readable count after the Gaussian estimate}.}

\teamanchor{proof.weighted-mollified-zero-test}%
For integral \(X\ge3\), take \(U=\lfloor X^{\Delta/4}\rfloor\) and
\(P=\lfloor X/U\rfloor\). Then \(N=UP\le X\) and \(P\ge U\).
For a primitive nonprincipal character \(\psi\) of conductor \(q\),
the P\'olya--Vinogradov inequality gives
\(|A_\psi(t)|\ll\sqrt q\log(2q)\) for \(A_\psi(t)=\sum_{n\le t}\psi(n)\).
Indeed its finite Fourier coefficients have modulus \(\sqrt q\) at
unit frequencies and vanish at nonunits. A change of variable shows
that the coefficients at units differ by unit-modulus factors.
For a nonunit frequency \(h\), multiplication by a unit congruent to
\(1\bmod q/(q,h)\) on which \(\psi\) is nontrivial fixes the additive
phase, giving zero. Such a unit exists by primitivity, and finite
Parseval then gives the magnitude \(\sqrt q\).
Fourier inversion and the geometric-sum bound
\(O(q/\min(h,q-h))\) for \(1\le h<q\) give the inequality.
For \(\Re s>0\), partial summation gives
\[
 L(s,\psi)-\sum_{n\le P}\psi(n)n^{-s}
   =-A_\psi(P)P^{-s}+s\int_P^\infty A_\psi(t)t^{-s-1}\,dt .
\]
At a tested zero \(\rho\) of \(L(s,\psi_d)\), we have \(q\le MX\),
\(\Re\rho\ge\beta_k\) and \(|\rho|\ll_k1\). The finite sum at \(\rho\)
therefore has size \(O_{M,k}(X^{1/2}P^{-\beta_k}\log(2X))\).
Multiply it by the finite sum
\(\mathcal M_U(s,\psi_d)=\sum_{m\le U}\mu(m)\psi_d(m)m^{-s}\),
where \(\mu\) is the M\"obius function.
Since \(|\mathcal M_U(\rho,\psi_d)|\ll_k U^{1-\beta_k}\) and
\(P\asymp X/U\), the product at \(\rho\) is
\[
 O_{M,k}\bigl(X^{-\Delta}U\log(2X)\bigr)
     =O_{M,k}\bigl(X^{-3\Delta/4}\log(2X)\bigr).
\]
The coefficient of \(n^{-s}\) in the product is
\(b_{U,P}(n)\psi_d(n)\), where
\[
 b_{U,P}(n)=\sum_{\substack{m\mid n\\m\le U,\ n/m\le P}}\mu(m).
\]
It satisfies \(b_{U,P}(1)=1\), \(b_{U,P}(n)=0\) for \(1<n\le U\),
and \(|b_{U,P}(n)|\le\tau(n)\), the number of divisors of \(n\).
Hence for all sufficiently large \(X\), uniformly over the tested zeros,
\begin{equation}\label{eq:weighted-mollified-zero-test}
 \left|E_d(\rho)\right|\ge\tfrac12,\qquad
 E_d(s)=\sum_{U<n\le N}b_{U,P}(n)\xi^*(n)\chi_d(n)n^{-s}.
\end{equation}

\teamanchor{proof.weighted-mollified-mean}%
Let \(S(Q_0)\) be the real primitive characters of conductor at most
\(Q_0\). Heath-Brown's quadratic large sieve for arbitrary complex
coefficients gives
\[
 \sum_{\chi\in S(Q_0)}
       \left|\sum_{n\le N}a_n\chi(n)\right|^2
 \ll_\alpha (Q_0N)^\alpha(Q_0+N)
     \sum_{\substack{n,m\le N\\nm\text{ a square}}}|a_n a_m|
 \qquad(\alpha>0).
\]
\cite[Corollary~2, p.~238]{f1heath-brown}\footnote{\href{https://matwbn.icm.edu.pl/ksiazki/aa/aa72/aa7234.pdf\#page=4}{Heath-Brown, original printing, Corollary~2 on p.~238}.}

Apply the inequality with \(Q_0=X\) to the distinct real primitive
characters \(\chi_d\). The fixed factor \(\xi^*(n)\) occurs in
\(a_n=b_{U,P}(n)\xi^*(n)n^{-s}\) on \(U<n\le N\).
For \(\Re s\ge1/2+\Delta/2\), the divisor bound
\(\tau(n)\ll_\Delta n^{\Delta/8}\) gives
\[
 |a_n|\ll_\Delta U^{-\Delta/4} n^{-1/2-\Delta/8}.
\]
Writing \(n=h r^2\), \(m=h t^2\) with a common squarefree \(h\) when
\(nm\) is a square shows that the square-pair sum is
\(O_\Delta(U^{-\Delta/2})\). The remaining sum is bounded by the
product of the three convergent sums of \(h^{-1-\Delta/4}\),
\(r^{-1-\Delta/4}\), and \(t^{-1-\Delta/4}\). Since \(N\le X\),
taking \(\alpha=\epsilon/2\) for \(\epsilon>0\) gives
\[
 \sum_{d\in\mathcal P(X)}|E_d(s)|^2
      \ll_{\Delta,\epsilon} X^{1+\epsilon} U^{-\Delta/2}
 \qquad (\Re s\ge1/2+\Delta/2).
\]

For each \(d\in\mathcal B_\xi(X)\), choose one zero \(\rho_d\) in its
box. The disk of radius \(\Delta/2\) about \(\rho_d\) lies in the fixed
rectangle
\[
 \mathcal R=\bigl\{s:1/2+\Delta/2\le\Re s\le1+\Delta/2,\
                          |\Im s|\le k+\Delta/2\bigr\}.
\]
The finite Dirichlet polynomial \(E_d\) is entire. The submean inequality
for \(|E_d|^2\) on that disk, followed by \eqref{eq:weighted-mollified-zero-test}, gives
\[
 \frac{\#\mathcal B_\xi(X)}4
 \le \frac{4}{\pi\Delta^2}
       \int_{\mathcal R}\sum_{d\in\mathcal P(X)}|E_d(s)|^2\,dA(s)
 \ll_{M,k,\epsilon} X^{1+\epsilon}U^{-\Delta/2}.
\]
Here \(dA\) is planar area. Since \(U\asymp X^{\Delta/4}\), take
\(\epsilon=\Delta^2/16\) to get the asserted count.

\teamanchor{proof.weighted-twist-choice}%
The union of the exceptional sets over all characters modulo \(M\)
has cardinality \(O_{M,k}(X^{1-1/(16k^2)})=o(X/\log X)\).
In the class \(a_0\), \eqref{eq:weighted-fixed-pnt} gives total von
Mangoldt weight asymptotic to \(X/\varphi(M)\), and the proper prime
powers contribute \(O(\sqrt X\log X)\). The class therefore has
\(\gg_M X/\log X\) primes up to \(X\). Choose one \(d_k\) outside the
union. This supplies \eqref{eq:weighted-twist-free} on the whole
character family before any first scale is chosen.

\teamanchor{proof.weighted-real-twist-density}%
For the real characters there is also a sharper estimate. For fixed
real \(\xi\), distinct \(d\) give distinct real primitive \(\psi_d\)
of conductor \(f_\xi d\le MX\). The
\hyperlink{proof.real-character-zero-density}{Heath-Brown zero-density
theorem} with conductor cutoff \(MX\), height \(k\), and left edge
\(\beta_k\) gives for the number of failing primes
\(O_{M,k,\alpha}(X^{1-4/(3k-2)+\alpha})\) for every \(\alpha>0\).
In particular \(\alpha=2/(3k-2)\) gives
\(O_{M,k}(X^{1-2/(3k-2)})\).
\cite[Corollary~1 and Theorem~3]{f1heath-brown}

\subsection*{One test square per block}

\teamanchor{proof.weighted-square-selection}%
Hold \(M,k,d_k\) fixed. For each integer \(j\ge1\), put \(X_j=2^{2^j}\).
Both endpoints of \([X_j,X_{j+1}]\) are squares, so its set of integer
squares is finite and nonempty. Choose the least square \(y_j\) in this
interval maximizing
\begin{equation}\label{eq:weighted-square-selection}
 S_{d_k,M}(y)-\frac{\log y}{\varphi(M)}-\tfrac12 y^{\beta_k-1}.
\end{equation}
This also maximizes \(R_{d_k,M}(y)-\tfrac12 y^{\beta_k-1}\), since the
centering constant \(c_{d_k,M}\) is common to all points. This is the
\hyperlink{proof.centered-ramp-maximum}{centered ramp maximum} on the
finite square set, with \(b=\lambda_{d_k,M}\), \(a=1/\varphi(M)\),
\(\alpha=\beta_k\) and \(\kappa=1/2\). The weights specify the test
squares before any arithmetic datum is chosen at them.

\teamanchor{proof.weighted-squares}%
At each selected \(y_j\) use the matching datum and row comparison with
\(d=d_k\) and \(\varepsilon=\varepsilon_k\). Put
\(\alpha_k=1/2+8\varepsilon_k=1/2+1/(2k)<\beta_k\). By
\eqref{eq:weighted-square-bound}, \(R_{d_k,M}(y_j)\le y_j^{\alpha_k-1}\)
for all large \(j\), since \(y_j\ge X_j\to\infty\) and the threshold is
fixed. Eventually \(y_j^{\alpha_k-1}\le\tfrac12 y_j^{\beta_k-1}\).
The centered ramp maximum rule gives
\(R_{d_k,M}(y)\le\tfrac12 y^{\beta_k-1}\) at every square in that interval.
The consecutive intervals \([X_j,X_{j+1}]\) cover the entire tail, hence
the last inequality holds at all sufficiently large squares.

For real \(x\ge1\) put \(y=\lceil\sqrt x\rceil^2\). Because the weights
are nonnegative and each kernel \((1/n-1/x)_+\) is nondecreasing in \(x\),
\(S_{d_k,M}(x)\le S_{d_k,M}(y)\). For sufficiently large \(x\) we obtain
\begin{align*}
 R_{d_k,M}(x)
 &\le R_{d_k,M}(y)+\frac{\log(y/x)}{\varphi(M)}\\
 &\le \tfrac12 x^{\beta_k-1}+\frac{2}{\varphi(M)\sqrt x}
 \le x^{\beta_k-1}.
\end{align*}
Here \(\beta_k<1\), and \(\log(y/x)\le2\log(1+x^{-1/2})\le2x^{-1/2}\).
The second half of the power bound absorbs the correction because
\(x^{\beta_k-1/2}\to\infty\). Thus
\begin{equation}\label{eq:weighted-real-bound}
 R_{d_k,M}(x)\le x^{\beta_k-1}\qquad\text{for every sufficiently large real }x.
\end{equation}
The size threshold may depend on \(k\) and on its chosen prime \(d_k\).
The \hyperlink{proof.weighted-window-variant}{power-window variant}
uses its own intervals with a positive-majorant criterion.

\subsection*{The polar transform and Landau's lemma}

\teamanchor{proof.weighted-polar}%
Apply \eqref{eq:common-ramp-mellin} with \(b=\lambda_{d,Q}\) and
\(a=1/\varphi(Q)\). The series identity \eqref{eq:weighted-series} gives
\(C_b=G_{d,Q}\) and \(c_{d,Q}=G_{d,Q}(1)\). Thus the function being
transformed is \(y\mathcal R_b(y)=yR_{d,Q}(y)\).
On \(\Re s>1\), set
\[
 \Phi_{d,Q}(s)=\int_1^\infty \bigl(yR_{d,Q}(y)\bigr)y^{-s-1}\,dy
 =\int_0^\infty R_{d,Q}(e^v)e^{(1-s)v}\,dv
\]
and the shared calculation gives
\begin{equation}\label{eq:weighted-polar}
 \Phi_{d,Q}(s)=\frac{F_{d,Q}(s)}{s(s-1)}
   -\frac{1}{\varphi(Q)(s-1)^2}-\frac{c_{d,Q}}{s-1}
   =\frac{G_{d,Q}(s)/s-c_{d,Q}}{s-1}.
\end{equation}
The formula continues meromorphically to \(\Re s>0\), and the shared
calculation removes the apparent pole at \(1\). The nonvanishing of the
regularized \(L\)-functions makes \(\Phi_{d,Q}\) holomorphic at every other
point with \(\Re s\ge1\). At each point
\(\rho\) with \(0<\Re\rho<1\), a logarithmic derivative has at most a
simple pole with residue equal to its zero order. Thus the poles there
are at most simple and the residue, taken to be zero at a regular point,
is
\begin{equation}\label{eq:weighted-residue}
 r_{d,Q}(\rho)=
 \frac{\displaystyle\sum_{\xi\bmod Q}
       \bigl(\operatorname{ord}_\rho\widetilde L_\xi-
             \operatorname{ord}_\rho\widetilde L_{\xi\chi_d}\bigr)}
      {\varphi(Q)\rho(1-\rho)}.
\end{equation}
This is the residue convention of the compact
\href{https://github.com/lsernau/grh-lean/blob/72fead102f663d5c7a407964a13dc0cf9d977207/grh-compact/src/New/Endgame/PrimeError.lean\#L66-L90}{polar-transform theorem}.

\teamanchor{proof.weighted-landau}%
Take \(d=d_k\), \(Q=M\), and put \(\beta=\beta_k\). For every real
\(\beta<\sigma<1\), all the twisted orders at \(\sigma\) vanish by
\eqref{eq:weighted-twist-free}, and the denominator in
\eqref{eq:weighted-residue} is positive. Thus \(r_{d,M}(\sigma)\ge0\).
At real \(\sigma\ge1\) there is no pole. These are the signs on the whole
real interval \(\sigma>\beta\) required by the compact
\href{https://github.com/lsernau/grh-lean/blob/72fead102f663d5c7a407964a13dc0cf9d977207/grh-compact/src/Lib/Analytic/Landau/PolarLandau.lean\#L84-L123}{one-sided Landau lemma}.

Put \(E(v)=R_{d,M}(e^v)\). The ramp kernels are continuous and locally only
finitely many are nonzero, so \(E\) is continuous. By
\eqref{eq:weighted-real-bound}, take
\(K\ge1\) large enough that \(g(v)=K e^{\beta v}-e^v E(v)\ge0\) for
all \(v\ge0\). Its Laplace transform is
\[
 P(s)=\int_0^\infty g(v)e^{-sv}\,dv
     =\frac{K}{s-\beta}-\Phi_{d,M}(s) \qquad (\Re s>2).
\]
Suppose the abscissa of convergence \(\sigma_c\) of the nonnegative
integral satisfies \(\sigma_c>\beta\). In \(\Re s>\sigma_c\) the integral
is holomorphic and agrees with the displayed expression by continuation.
At \(\sigma_c\) that expression has at most a simple pole, with residue
\(-r_{d,M}(\sigma_c)\le0\). For real \(s>\sigma_c\), the integral gives
\((s-\sigma_c)P(s)\ge0\). Its limit as \(s\downarrow\sigma_c\) is the
residue, hence the residue is zero and \(P\) is holomorphic in a
neighborhood of \(\sigma_c\).

Choose a real center \(a>\sigma_c\) so close to \(\sigma_c\) that the
Taylor series of \(P\) at \(a\) converges at some real \(a-\delta<\sigma_c\).
The derivatives at \(a\) are obtained by differentiating the integral in
its half-plane of convergence. Positivity of \(g\) and monotone
convergence give
\[
 \sum_{j\ge0}\frac{P^{(j)}(a)(-\delta)^j}{j!}
 =\int_0^\infty g(v)e^{-av}\sum_{j\ge0}\frac{(\delta v)^j}{j!}\,dv
 =\int_0^\infty g(v)e^{-(a-\delta)v}\,dv.
\]
The left side is finite, contradicting the abscissa. It follows that
\(P\) is holomorphic for \(\Re s>\beta\), and by the continuation formula
so is \(\Phi_{d,M}\). In particular \(r_{d,M}(\rho)=0\) for
\(\Re\rho>\beta\).

\subsection*{The classical zeta target}

\teamanchor{proof.weighted-rh}%
Suppose that \(\zeta(\rho)=0\) with \(1/2<\Re\rho<1\). Choose
\(k\ge13\) with \(\beta_k<\Re\rho\) and \(|\Im\rho|\le k\), and use the
same \(d_k\) in the arithmetic bound and in the residue formula. The
conductor-one principal character induces the principal character
\(\xi_0\bmod M\). By \eqref{eq:common-character-conductor},
\[
 \widetilde L_{\xi_0}(s)=(s-1)\zeta(s)\prod_{p\mid M}(1-p^{-s}).
\]
The finite factors are nonzero at \(\rho\), so the
untwisted principal term contributes a positive zero order to
the numerator of \eqref{eq:weighted-residue}. All other untwisted orders
are nonnegative and all twisted orders vanish by
\eqref{eq:weighted-twist-free}. The residue is then nonzero, contrary to
Landau's lemma. This proves the right-strip exclusion for the classical
zeta function.

The compact source's
\href{https://github.com/lsernau/grh-lean/blob/72fead102f663d5c7a407964a13dc0cf9d977207/grh-compact/src/New/Main.lean\#L169-L197}{Dirichlet GRH deduction}
uses the arithmetic at modulus \(18480q\) for each target modulus \(q\).
Here the full auxiliary character family at the single modulus \(18480\)
enters the residue test for the classical target.

\subsection*{From the right half of the strip to classical RH}
\teamanchor{proof.zeta-reflection}%
Suppose that \(\zeta(s)\ne0\) whenever \(1/2<\Re s<1\).
Then every zero of \(\zeta\) other than the negative even integers
has real part \(1/2\). Both routes in this paper use this implication.

The \hyperlink{proof.dirichlet-zero-free-line}{modulus-one boundary
calculation} gives nonvanishing on \(\Re s\ge1\) away from the
simple pole at \(1\).
For \(\Re s<1/2\) and \(s\ne0\), the functional equation gives
\[
  \zeta(s)=2^s\pi^{s-1}\sin(\pi s/2)\Gamma(1-s)\zeta(1-s).
\]
Since \(\Re(1-s)>1/2\) and \(1-s\ne1\), the hypothesis and the
boundary result make \(\zeta(1-s)\) finite and nonzero.
The gamma factor is finite and nonzero in this region.
The sine factor vanishes there precisely at the negative even integers,
and the remaining point \(s=0\) has \(\zeta(0)=-1/2\).
\cite[(25.4.2), (25.6.1) and \S5.2(i)]{dlmf}
This proves the stated implication.

